\documentclass[11pt,a4paper]{article}
\usepackage[margin=25mm]{geometry}
\usepackage{amsmath,amssymb,amsthm,booktabs,array,xcolor,fancyhdr,longtable}
\usepackage[T1]{fontenc}
\usepackage{lmodern}
\usepackage[hidelinks]{hyperref}
\setlength{\emergencystretch}{3em}
\pagestyle{fancy}\fancyhf{}
\fancyhead[L]{\small The iid Berry--Esseen constant}
\fancyhead[R]{\small Computer-assisted argument}
\fancyfoot[C]{\thepage}
\newtheorem{theorem}{Theorem}[section]
\newtheorem{lemma}[theorem]{Lemma}
\newtheorem{proposition}[theorem]{Proposition}
\newtheorem{definition}[theorem]{Definition}
\newcommand{\E}{\mathbb E}
\newcommand{\Pp}{\mathbb P}
\newcommand{\R}{\mathbb R}
\newcommand{\CE}{C_{\mathrm E}}
\newcommand{\supp}{\operatorname{supp}}
\title{The iid Berry--Esseen constant\\
\large A computer-assisted argument with exact certificates}
\author{Working preprint}
\date{27 October 2026}
\begin{document}
\maketitle
\begin{abstract}
We assemble an argument for the optimal iid Berry--Esseen constant from
ordinary analytical reductions and finite exact certificates. The proof
selects an unrestricted extremal law, preserves that law through contact,
resource and conditional-count estimates, and covers all positive integers
by thirteen compatible contributions. The infinite tail has its own
three-loss base and least-failure argument. We give the classical lower
approximation separately, prove the binomial normalization extension across
one half, and provide a companion source supplement with the longer
analytical derivations and their logical dependencies.
\end{abstract}

\section{Scope and target}
Let \(\mathcal P\) be the real probability laws with mean zero, variance one,
and finite third absolute moment \(\rho(P)=\E_P|X|^3\). For iid variables
with law \(P\), let \(S_n=X_1+\cdots+X_n\) and
\[
 \Delta_n(P)=\sup_{H\in\R}
       |\Pp_P(S_n\le H)-\Phi(H/\sqrt n)|,\qquad
 R_n(P)=\frac{\sqrt n\,\Delta_n(P)}{\rho(P)}.
\]
Set
\[
 C_n=\sup_{P\in\mathcal P}R_n(P),\qquad C^*=\sup_{n\ge1}C_n,\qquad
 \CE=\frac{3+\sqrt{10}}{6\sqrt{2\pi}}.
\]
\begin{theorem}[Conclusion from the exact certificates]\label{thm:global}
Assume that the complete source-pinned numerical predicates of the thirteen
families in Table~\ref{tab:roles} are true at their exact scopes. With the
published inputs specified below and the ordinary analytical reductions
proved here and in the precise source supplement~\cite{source},
\[
 C_n\le\CE\quad(n\ge1),\qquad C^*=\CE.
\]
\end{theorem}
The thirteen source families have been assembled and the scoped independent
mathematical review is complete. The main argument below explains how their
exact numerical conclusions enter ordinary mathematical implications.
The companion \path{analytic-supplement.txt} contains the longer proofs,
with original labels and a lemma-to-main-claim index. Verification status is
summarized once in Section~\ref{sec:open}.

Affine standardization needs no one-summand bound. For a law $Z$ with mean
$\mu$, variance $\sigma^2>0$ and third central absolute moment $\beta_3$,
set $X=(Z-\mu)/\sigma$. Then $\rho=\beta_3/\sigma^3$ and
\[
 \Pp\!\left(\frac{\sum Z_i-n\mu}{\sigma\sqrt n}\le z\right)
 =\Pp\!\left(\sum X_i\le\sqrt n\,z\right).
\]
The result therefore has its usual arbitrary-mean/positive-variance form.

Throughout, \(H\) is a \emph{raw sum} threshold. Its normalized version is
\(H/\sqrt n\). All terminating decimals used by a checker mean exact
rationals. A numerical relaxation may contain nonrealizable moment points:
soundness requires that it contain every actual law, not that every point
be a law.

\section{The independent classical lower approximation}
Set $p=(4-\sqrt{10})/2$, $q=1-p$, and $X=(B-p)/\sqrt{pq}$ for a
Bernoulli variable $B$. Then $\E X=0$, $\E X^2=1$ and
$\rho=(p^2+q^2)/\sqrt{pq}$. The fixed-law limit in
Schulz~\cite[Remark~4.1, printed p.~15]{schulz} gives, in the present normalization,
$R_n(X)\to\CE$: exact algebra gives $p^2+q^2=10-3\sqrt{10}$ and
$(3+q-p)/(p^2+q^2)=3+\sqrt{10}$. The threshold substitutions are
\[
 H=\sqrt n\,z,\qquad x=np+\sqrt{pq}\,H=np+\sqrt{npq}\,z.
\]
Each ranges over the whole real line. For every $b<\CE$, choose
$\varepsilon=(\CE-b)/2$ and a positive integer $n$ with
$R_n(X)>\CE-\varepsilon>b$. Thus $C^*\ge\CE$, independently of the upper
proof and without finite-index attainment. The all-parameter finite-binomial
upper theorem in~\cite[Theorem~1]{schulz} is a separate input, used on the
exact zero-noise branch below; neither theorem is proved by a PDF hash.

\section{The unrestricted owner and its contact}
\begin{lemma}[Fixed-index attainment]\label{lem:attain}
For each fixed \(n\ge1\), \(0<C_n\le\sqrt n\), and some law \(P\in\mathcal P\)
and threshold \(H\) attain \(C_n\) as a positive closed discrepancy:
\[
 \Pp_P(S_n\le H)-\Phi(H/\sqrt n)=\Delta_n(P),\qquad R_n(P)=C_n.
\]
\end{lemma}
\begin{proof}
A negative closed discrepancy reflects into a positive strict discrepancy.
Approaching that strict threshold from below, continuity of \(\Phi\) gives
the same supremum with positive closed discrepancies. Since \(\rho\ge1\),
the upper bound is immediate. A symmetric two-point law, evaluated at the
upper endpoint of its sum, gives positivity.

Take positive scores \(r_j\to C_n\), with \(r_j\ge c>0\). Then
\(\rho_j\le\sqrt n/c=:M\). Tightness gives a subsequence \(P_j\Rightarrow P\);
also take \(\rho_j\to b\). The estimates
\[
 \int_{|x|>A}x^2\,dP_j\le M/A,\qquad
 \int_{|x|>A}|x|\,dP_j\le M/A^2
\]
preserve the first two moments. Lower semicontinuity gives
\(1\le\rho(P)\le b<\infty\). Chebyshev's inequality for negative thresholds,
and the Gaussian upper tail for positive thresholds, imply that the chosen
raw thresholds satisfy \(|H_j|\le n^{3/4}/\sqrt c\). Pass to \(H_j\to H\).
For fixed \(n\), weak convergence of products and the continuous sum map
give \(P_j^{*n}\Rightarrow P^{*n}\). Closed-set Portmanteau at \(H+\varepsilon\),
followed by \(\varepsilon\downarrow0\), yields
\[
 F_{n,P}(H)-\Phi(H/\sqrt n)\ge C_n b/\sqrt n.
\]
Consequently \(C_n\ge R_n(P)\ge C_n b/\rho(P)\ge C_n\).
All inequalities are equalities. In particular \(\rho(P)=b\); third-moment
convergence is a consequence of extremality, not an assumed uniform
integrability of third moments. The proof includes \(C_n=\CE\).
\end{proof}

The law used below is the owner \emph{selected by this lemma}, after
orientation. We do not begin with an arbitrary law and pretend it remains
an optimizer after extraction. Once selected, the law and raw threshold
are fixed through every subsequent support, resource and count statement.

For \(n\ge2\), write \(C=C_n\), \(F=F_{n,P}(H)\),
\(z=H/\sqrt n\), \(\gamma=\E X^3\), and \(\beta=\E[X|X|]\).
The exact contact statement used by the proof is
\begin{align}
 \rho I(y)={}&n^{3/2}\{F_{n-1,P}(H-y)-F\}
      +ny\phi(z)+\tfrac12H\phi(z)(y^2-1)\notag\\
 &+\tfrac32C\rho(y^2-1)
      -C\{|y|^3-\rho-3y\beta\},                  \label{eq:contact}\\
 I(y)&\le0\quad(y\in\R),\qquad
 I(y)=0\quad(y\in\supp P).                     \label{eq:contact-zero}
\end{align}
\begin{proof}[Contact calculation]
Contaminate the raw law by $P_\varepsilon=(1-\varepsilon)P+
\varepsilon\delta_y$ and standardize. Its mean and variance derivatives
are $y$ and $y^2-1$. Its centered third-absolute-moment derivative is
$|y|^3-\rho-3y\beta$, justified by an integrable quadratic derivative bound.
The raw closed-convolution derivative is
$n\{F_{n-1,P}(H-y)-F\}$, by the finite mixture expansion; the normal
argument derivative is $-\sqrt n\,y-H(y^2-1)/(2\sqrt n)$.
Differentiating the normalized score gives exactly $\rho I(y)$;
extremality makes it nonpositive. Its expectation is zero. The contact
majorant is continuous while $F_{n-1,P}(H-y)$ is upper semicontinuous;
thus strict negativity at a support point would give a negative integral
on a neighborhood. This proves equality on the whole closed support.
The cubic leading majorant confines that support before higher moments
are introduced. See~\cite[lines410--463]{source} for the full atom-safe
variation and compactness details.
\end{proof}

\section{The same-law predecessor step}\label{sec:predecessor}
For a standardized law \(P\), define
\[
 q_n(P,C)=n\max\{R_{n-1}(P)-C,0\}.
\]
\begin{lemma}[Weak predecessor is sufficient]\label{lem:qzero}
If \(C_{n-1}\le\CE\), then for every \(P\in\mathcal P\) and every
\(C\ge\CE\), \(q_n(P,C)=0\).
\end{lemma}
\begin{proof}
The definition of the unrestricted supremum gives
\(R_{n-1}(P)\le C_{n-1}\le\CE\le C\). The positive part is zero.
\end{proof}
For the selected owner, use \(C=R_n(P)=C_n\). Thus every previously proved
implication using an allowance \(q_n(P,C_n)\le Q\), \(Q\ge0\), remains valid.
No contact at \(n-1\), strict predecessor, or monotonicity in \(n\) is inferred.
Own-index score caps and resource bounds are separate premises and are not
replaced by this lemma.

The source manuscript proves exactly this step at
\texttt{indcore:lem:replacement}, lines 4777--4789. Some older resource
descriptions assume the stronger \(C_{n-1}<\CE\). The assembly uses the
actual weaker allowance hypothesis of their support/resource lemmas, together
with Lemma~\ref{lem:qzero}; it does not assert that stronger predecessor.

\section{Support and independently licensed Fourier bounds}
Put \(s=\sqrt{n/(n-1)}\), \(e=n(s-1)-1/2\), and
\(A=R_{n-1}(P)\). From contact one obtains, at a support point \(t\),
\[
 0\le\mathcal H_n(H,t)
   +C\{-|t|^3+\tfrac32\rho t^2+3\beta t+\rho e\}
   +\rho s\,n(A-C).
\]
Here
\[
 \mathcal H_n(H,t)=n\sqrt n\{
 \Phi((H-t)/\sqrt{n-1})-\Phi(H/\sqrt n)\}
 +\{nt+\tfrac12H(t^2-1)\}\phi(H/\sqrt n).
\]
For bounds \(\rho\le r\), \(|\beta|\le b\), \(s\le S\), \(e\le E\),
\(0<c\le C\), and \(q_n(P,C)\le Q\), define
\[
 Q_0(v)=-v^3+\tfrac32rv^2+3bv+rE,\qquad
 L(v)=\tfrac25(v^2/2+1/2+E)+cQ_0'(v).
\]
The direct support lemma in~\cite{source}, lines 492 onward, requires
\[
 Q_0(R),Q_0'(R),Q_0''(R),L(R),L'(R)<0,\quad 1/5-3c<0,
\]
and, for both signs \(\eta\),
\[
 \sup_{H\in\R}\mathcal H_n(H,\eta R)+cQ_0(R)+rSQ<0.
\]
Since \(Q_0'''=-6\) and \(L''=2/5-6c<0\), these inequalities propagate
along both complete closed rays \(v\ge R\). They contradict contact on
those rays, proving \(\supp P\subset(-R,R)\). Whole-threshold Gaussian
enclosures, including both infinite threshold tails, are required at the
endpoint test. Sampling thresholds does not prove that supremum bound.

Fourier bounds serve different roles. An unrestricted auxiliary ceiling
may license an own-score cap before support is known. A supported Fourier
ceiling applies only after that same law's support has been acquired.
Their normalization and endpoint terms must remain exact. In the
20--24 family the source Prawitz reduction uses the doubled kernel
\[
 K(t)=(1-t)+i\{(1-t)\cot(\pi t)+1/\pi\}.
\]
The four smoothing terms, the omitted endpoint pieces, and the infinite
Gaussian tail are part of the certificate. A finite-window integral alone
is not a Kolmogorov bound. The independent scalar-cosine certificate and
Prawitz inequality are explicit inputs; numerical authentication cannot
grant these analytic statements.

The exact iid literature inputs are~\cite{shevtsova2013}
\[
 R_n(P)\le\min\{.4690,\ .3322(1+.429/\rho),\ .3031(1+.646/\rho)\}.
\]
They apply for every $n\ge1$: the source normalization for $X_i/\sqrt n$
has $\ell_n=\rho/\sqrt n$ and $\tau_n=1/\sqrt n$.
The n-generic compiled complement deliberately weakens $.3322$ to $.3328$;
this also gives a valid bound, decreasing below $.4097$ at $\rho=19/10$.
The smoothing convention is fixed by~\cite[Lemma~4.1, printed p.~263]{shevtsova2012}.
The directed inequality $.4097<\CE$ is separate from the coarse
$.4095<\CE$ used for cap/excess estimates.

\section{One marginal and the conditional-count identity}
Keep the same \(P\). Set
\[
 p=\Pp(X>0),\quad q=1-p,\quad v=pq,\quad
 m=\E X_+=\E X_->0,\quad h=v/m.
\]
The mean-zero, nondegenerate law has \(0<p<1\). With
\(I=\mathbf1_{\{X>0\}}\), write
\[
 hX=I-p+U,\qquad \E(U\mid I)=0.
\]
A zero atom belongs to the nonpositive class. Let \(U_i\) be the two
conditional laws and \(w_i=\E U_i^2\). Then
\[
 s=qw_0+pw_1=h^2-v,\qquad \lambda=ns.
\]
In standardized physical units $\theta=s/(v+s)$, since $h^2=v+s$.
Consequently $\lambda=npq\theta/(1-\theta)$.
If \(s>0\), let \(\alpha=pw_1/s\); if \(s=0\), both conditional residuals
are zero and no division by \(s\) is made.

Conditional on \(j\) positive draws, the sum of residuals has law
\[
 W_j=U_{0,1}+\cdots+U_{0,n-j}+U_{1,1}+\cdots+U_{1,j},
\]
with independent summands. For \(b_j=\binom njp^jq^{n-j}\) and
\(np+hH=k+u\),
\begin{align}
 \Pp(S_n\le H)
 &=\sum_{j=0}^n b_j\Pp(W_j\le k+u-j)\notag\\
 &\le F_{\mathrm{Bin}(n,p)}(k)-b_k\Pp(W_k>u)
       +\sum_{j>k}b_j\Pp(W_j\le u-(j-k)).       \label{eq:count}
\end{align}
Indeed \(hS_n=K-np+\sum U_i\); condition on \(K\), retain \(j=k\)
exactly, and bound \(j<k\) by its full mass. The outgoing event is strict
because its complement is the closed event \(\{W_k\le u\}\).
Every incoming event remains closed. This identity works on both
probability halves without reflecting \(p,\beta,\gamma\).

\subsection*{Joint resources and units}
Suppose the actual central support gap is \((-A,B)\). The scaled residual
endpoints are
\[
 a_0=p-hA,\quad a_1=q-hB,\qquad
 U_0\in[-O_0,a_0],\quad U_1\in[-a_1,O_1].
\]
The raw inward displacement is
\[
 d_{\rm raw}=(m/q-A)+(m/p-B),\qquad a_0+a_1=h\,d_{\rm raw}.
\]
Thus raw inward displacement and its scaled COUNT coordinate cannot be
identified without their factor \(h\). In these scaled coordinates, a
licensed joint plane has the form
\[
 \frac{a_0+a_1}{d_0}+\frac{\lambda}{\Lambda}\le1.
\]
A later, smaller marginal cap for \(\lambda\) does not replace the
intercept \(\Lambda\) in this plane. A zero resource budget has its own
deterministic conclusion. The physical defect, its affine halfplane and
the two signed quadratic constraints must all concern this same law.
Their full proofs are the source labels
\texttt{rescore:joint-resource}, \texttt{countcore:physicalD} and
\texttt{countcore:Dhalfplane}; they are distinct ordinary analytical inputs,
with proofs in that supplement.

\section{What a finite COUNT certificate proves}
On a complete acquired parameter box, the checker combines valid outgoing
lower and incoming upper probability bounds in~\eqref{eq:count}, then
subtracts a lower Gaussian bound at the \emph{same} phase. If \(B_k\),
\(b_k^-\), \(b_j^+\), \(L_k\), \(A_j\), \(N\) denote those bounds, it
certifies a whole-box upper for
\[
 B_k-b_k^-L_k+\sum_{j>k}b_j^+A_j-N.
\]
This combined expression retains phase cancellation. Separate maxima at
unrelated phases would not be the same proof. The positive error is then
multiplied by an upper enclosure of \(\sqrt n\) and divided by a positive
lower enclosure of the same \(\rho\). Every closed normalization piece is
included. Convex vertex maximization is performed before minimizing over
a finite menu of strategies; a pointwise minimum is not assumed convex.

If the acquired bounds imply
\[
 np_- -h_+H_-\ge-1/2,\qquad np_+ +h_+H_+\le n+1/2,
\]
every threshold has a nearest count \(k=0,\ldots,n\) and
\(u\in[-1/2,1/2]\). Otherwise the remaining thresholds require separate
proved tail bounds. The count notation does not discard them.

Each of the \(n+1\) original product roots is covered by a finite binary
tree. A strict interior rational cut retains \emph{both closed children}.
Every leaf either proves a necessary-condition contradiction or a
whole-box normalized score upper below a rational \(c<\CE\).
Induction up the finite tree proves its root statement; union over all
roots proves the complete comparison, including ties, boundaries and
zero-noise faces. Missing roots, unsuccessful guards and unfinished
searches are not exclusion certificates. This is the precise finite
closed-cover implication of~\cite{source},
\texttt{countcore:scoretheorem} and \texttt{countcore:forest}.

The entire matching contract and physical resource-to-COUNT implication
remain mandatory. A local COUNT acceptance does not acquire support,
an own-index cap, contact, or a new all-law theorem by itself.

\section{The literal 20--24 chain}
This completed source block makes the assembly concrete.
The full numerical role set has 186 entries:
163 ordinary forests, eight INITIAL/resource cores, three Fourier packets,
one same-law 19/20 feedback packet, one high-support packet, the shared
scalar certificate, one full actual20 conjunction and eight full band joins.
Already checked inputs are reused with their existing provenance; this
list is not a request to repeat their numerical work.

For each \(n=21,\ldots,24\), the closed low-rho bands and initial support
radii are
\[
 \begin{array}{c|c|c}
 &\rho\text{ interval}&\text{radius}\\ \hline
 B0&[1,11/10]&23/10\\
 BN&[11/10,9/8]&12/5 .
 \end{array}
\]
Every own cap \(419/1000\) comes from a separate unrestricted auxiliary
Fourier bound at the \emph{actual} index. Each core starts without an
acquired defect seed, runs eight feedback covers and a ninth export,
and accounts for all 2592 signed cells per cover and all 100 closed
probability panels. Its quantitative resource meaning is required even
if a discovery target happened to be reached earlier.

Actual20 uses 49 complete forests, with its original whole-panel
two/four-piece unions. The eight later bands use 114 forests and 800
complete probability dispositions. All 1000 panels in this five-index
block are accounted for by physical contradiction, strict scalar
contradiction, or the full matching COUNT. Every probability seam is
closed. The high-rho complement uses supported Fourier on \([9/8,6/5]\),
support-free Fourier on \([6/5,19/10]\), and the full published infinite
tail. Radius-3 support for that supported branch is a separate premise.

\begin{proposition}[Conditional five-index composition]
Suppose the actual all-law \(C_{19}\le\CE\) theorem is supplied, every
numerical input above passes its complete mathematical interpreter, and
the exact attainment/contact, Fourier, support, resource, realization and
cover implications are proved or accepted at their stated scope. Then
\(C_n\le\CE\) for \(n=20,\ldots,24\).
\end{proposition}
\begin{proof}
At the next literal index assume \(C_n>\CE\), choose the unrestricted owner
by Lemma~\ref{lem:attain}, and freeze that law and raw threshold.
Lemma~\ref{lem:qzero} licenses the same-law predecessor allowance.
The separate actual-index own caps, full rays, INITIAL and resource
implications place the owner in the complete low/high cover.
Every part contradicts its score or a necessary physical condition.
Thus \(C_n\le\CE\). Repeat at the next integer. No desired next bound,
sample-size monotonicity, or contact at a transferred index is an input.
\end{proof}

The current candidate supplies the actual whole41 predicate, completed
cached185 meanings, genuine incoming19 theorem and all12 analytical
implications. The large numerical computations remain qualified actual
source-execution premises, rather than a status-field change.

\section{The thirteen roles and the unbounded tail}
\begin{longtable}{@{}p{.10\linewidth}p{.48\linewidth}p{.34\linewidth}@{}}
\caption{Exact upper partition. Every finite index1--48 occurs once.}
\label{tab:roles}\\\toprule Index&Mathematical role&Complete source/numerical meaning\\\midrule
\endfirsthead\toprule Index&Mathematical role&Complete source/numerical meaning\\\midrule\endhead
1&$C_1\le2/5<\CE$; Appendix~\ref{app:n1}.&150 closed cubic-dual cells, independent exact checks.\\
2&$C_2\le2\sqrt2/7<\CE$, full unrestricted tensor reduction.&Two independent complete polynomial checkers.\\
3&$C_3\le59\sqrt3/250<\CE$, 59-feature null polynomial.&Full scalar, identity and cube predicates.\\
4&$C_4\le.4097<\CE$ from quantitative G3.&Both ROOT4 whole260-parent passes, HIGH/LOW and reserves.\\
5&New one-free G3$\to$G4$\to$same-owner $q_5=0$ contradiction.&Whole region/resource/COUNT cover;13 retained rules, no $.3704$ premise.\\
6--9&Exact source implication with the new proof of the identical weak C5 parent.&Whole original contexts, explicit parent substitution.\\
10--19&Exact parent substitution from the genuine6--9 theorem.&Whole original resource/COUNT contexts and closed subdomains.\\
20--24&Increasing weak induction, actual own caps $<.419$.&186 meanings,12 analytical rules,163 COUNTs,1000 panels.\\
25&New low-only $[1,6/5]$ cover and independent unrestricted high-rho branch.&Six original calls,52 complete COUNT interfaces, zero noise.\\
26--39&New low-only weak induction; actual support35 and separate positive39 cap.&300 COUNTs, whole7 partitions,35 support predicates, physical interfaces.\\
40--42&New floor40 theorem for actual $n\in\{40,41,42\}$.&CORE304+2/9/INITIAL2; honest original8+new2 acquisition conjunction.\\
43--48&Independent positive ceilings and finite least failure, no sharp42 premise.&Original183 complete parents; distinct floor42 LOW family.\\
$n\ge49$&Independent positive ceilings, strong46, uniform local exclusion.&14 top-level plus41 retained execution predicates; full closed cover and rays.\\\bottomrule
\end{longtable}
The principal finite dependency is G3$\to$G4$\to$C5$\to$C6--9$\to$C10--19
$\to$C20--24$\to$C25$\to$C26--39$\to$C40--42. N1 and G2 are independent.
The unrestricted high-rho25--48 theorem is independent. The43--48 and49-tail
roles have their own positive/coarse prerequisites: no sharp42$\to$43 edge
and no sharp1--48 premise in the tail are used. Proof-parent substitution
requires an identical statement and retains the complete numerical context.

\subsection*{Weighted floors and signed corrections}
A floor row $N$ is applied only with $nD_{\rm actual}\le Nd$ and both
the defect ceiling and noise price scaled by $N/n$; see Appendix~\ref{app:resource}.
All signed moment, support, probability and threshold fields must still match.
For real $I$ and $40\le n\le42$,
$I/n\le\max(I/40,I/42)$. For a signed bracket $B$ and multiplier in
$[a,b]$, its lower product is $\min(aB,bB)$. These are affine endpoint
inequalities, requiring no positive sign for $I$ or $B$.
The paired portfolio retains $M<3$, $g^2<2$, $R>2g$ and its positive
damping gates. Negative skewness uses the upper density endpoint, or its
explicit equivalent correction, before completing squares. Both raw and
barrier-subtracted curvatures are corrected. The ordinary active-override
proof, licenses this reuse.

\subsection*{Three losses at49 and the least later failure}
The independent auxiliary branch gives
\[
 C_j\le\CE+\frac{.7}{j}\ (45\le j\le50),\qquad
 C_{46}<\CE+\frac{.54}{46}.
\]
Write $d_n=n(C_{n-1}-C_n)_+$. Its uniform local theorem excludes a positive contrary owner for every
actual $n\ge47$ with excess at most $.7/n$ and $d_n\le.19$, current cap $C_n<.4247$, and deleted caps
$C_{n-1},C_{n-2}<.4253$. The entry theorem acquires the same-law support,
probability, moment and resource hypotheses independently of sharpness.
Conditional-count, signed-deletion, confinement, marked-crossing and
scaled-small-noise arguments prove $\rho_B(C-C_B)<-\lambda/1000$ for
every $\lambda>0$. At $\lambda=0$, centering gives $U=0$ almost surely
and Schulz's separate finite-binomial theorem gives $C_B<\CE$.
The full ordinary proofs are~\cite[lines7647--8064 and8217--8887]{source};
their complete numerical predicates, signed gates and closed seams are
required, not inferred from a PASS or finite sampling.

If $C_{49}>\CE$, the rational caps license the local theorem successively
at49,48,47. Each failure forces $C_{n-1}>C_n+.19/n$. Hence
\[
 C_{46}>\CE+\frac{19}{100}\left(\frac1{47}+\frac1{48}+\frac1{49}\right)
 >\CE+\frac{27}{50\cdot46},
\]
since the excess difference is exactly $35419/254251200>0$.
This proves weak sharpness at49 without sharp1--48. Each application may
choose its own owner; its internal same-law premises never change.

For the least later violation $n\ge50$, the previous weak sharp theorem
gives zero same-law decrement. The support-free smoothing bound is
\[
 C_n\le\sqrt{\frac n{n-1}}C_{n-1}
 +\frac{\sqrt n}{15(n-1)^{3/2}}+\frac2{25}\frac{\sqrt n}{(n-1)^2}.
\]
It follows by conditioning on one summand and using third-absolute-moment
Gaussian Taylor and variance heat remainders before either supremum
\cite[lines7463--7506]{source}; no support or higher moment is assumed.
For $n\ge50$, $\CE<409733/10^6$ and $\sqrt{n/(n-1)}<1011/1000$ give
\[
 n(C_n-\CE)<\frac{409733}{10^6}\frac{25}{49}
 +\frac{50}{49}\frac{1011}{15000}
 +\frac2{25}\frac{50}{49}\frac{1011}{7000}<.29.
\]
Thus the local theorem applies again. At $n=50$ its second predecessor48
uses the independent coarse bound $C_{48}\le\CE+.7/48<.4253$; thereafter
both deleted indices are earlier sharp ones. Contradiction proves the
entire unbounded tail. The thirteen-part cover and the separate lower
approximation prove Theorem~\ref{thm:global} under its stated predicates.

\section{Verification status}\label{sec:open}
The original source computations, global mathematical composition and scoped
independent analytical/checker-source review are complete. The numerical
hypotheses in Theorem~\ref{thm:global} are supplied by qualified complete
source-bound executions. A full fresh public end-to-end run is still in
progress; this version does not report it as completed or claim
\texttt{full\_pass}. Exact arithmetic, whole-domain native acceptance and
source/runtime/parent bindings are retained in the reproducibility package.

The accompanying ordinary proofs establish the analytical implications;
execution records identify the numerical computations. Complete Lean
formalization of the analysis and checker semantics remains future work;
the existing assembly/limit bridges verify conditional implications.
AI assistance was used in development, source review and preparation of
this working preprint.

\appendix
\section{The one-summand dual proof}\label{app:n1}
For each rational cell $[a,b]$ of a closed cover of $[-2,1]$, the finite
input supplies $u,v,L_a\in\mathbb Q$, $L_a\le\Phi(a)$, with
\[
 L_a-v+ux+vx^2+\tfrac25|x|^3-\mathbf1_{\{x\le b\}}\ge0\quad(x\in\R).
\]
Whole-line polynomial/Sturm checks include both infinite rays.
Expectation cancels $-v+v\E X^2$ and $u\E X$, giving
$F_X(b)-L_a\le(2/5)\rho$. Monotonicity gives the same upper bound for
$F_X(t)-\Phi(t)$ on $[a,b]$. At $t\le-2$, Cantelli gives $F_X(t)\le1/5$;
at $t\ge1$, $1-\Phi(1)<2/5$ (for example $\phi(0)>3/8$ and
$e^{-s^2/2}\ge1-s^2/2$ give $\Phi(1)>13/16$). Since $\rho\ge1$,
these suffice. Finally
$\Phi(t)-F_X(t)\le F_{-X}(-t)-\Phi(-t)$, since
$F_{-X}(-t)=1-F_X(t-)$, so reflection includes atoms.
The150-cell certificate and independent exact checks are finite inputs;
this ordinary implication needs no general-law $.3704$ claim.

\section{Physical defect and weighted units}\label{app:resource}
Put $\sigma=\sqrt{10}$, $\tau(t)=\E[X\mathbf1_{\{X\ge t\}}]$ for the
bounded owner. Then $\tau\ge0$, $\int\tau=1$ and $2\int|t|\tau(t)dt=\rho$.
For complementary gaps $(a,b)$ of its closed support let $\tau_{a,b}$ be
the constant interior value, and put
\[
 J=\sum_{(a,b)}\tau_{a,b}(b-a)^2,\qquad D=(3+\sigma)\rho-\gamma-3J.
\]
The sums converge by Tonelli and bounded support. If zero is in the support,
$(b-a)^2\le2\int_a^b|t|dt$ gives $J\le\rho$ and $D\ge0$. Otherwise let
$(-A,B)$ be the central gap, $m=\E X_+=\E X_-$,
$T_+=\E X_+^3-mB^2$, $T_-=\E X_-^3-mA^2$, $J_o=J-m(A+B)^2$.
On positive outer gaps use $(b-a)^2\le2\int_a^b(t-B)dt$ and reflect;
thus $J_o\le T_++T_-$. Direct expansion gives
\[
 D=m(\sqrt{\sigma+1}A-\sqrt{\sigma-1}B)^2
 +(\sigma-1)T_++(\sigma+1)T_-+3(T_++T_--J_o)\ge0.
\]
For the actual conditional means $\bar A=m/q$, $\bar B=m/p$, write
$\Delta_A=\bar A-A$, $\Delta_B=\bar B-B$ and let $\theta_+,\theta_-$
be the positive/negative conditional centered second moments with class
weights included; $\theta=\theta_++\theta_-$. Conditional cubic expansion
and the stronger outer estimates give the simultaneous ordinary budget
\begin{align*}
 D\ge{}&\frac{(6+2\sigma)(1-\theta)^{3/2}}{\sqrt{pq}}(p-p_*)^2
 +6m\{(A+B)(\Delta_A+\Delta_B)+\Delta_A\Delta_B\}\\
 &+3(\sigma+1)B\theta_++3(\sigma+3)A\theta_-,\qquad p_*=(4-\sigma)/2.
\end{align*}
Indeed the positive cubic residual is at least $3B\theta_+$ and its
reflection handles the other sign. The deterministic remainder is exactly
the second term, and
$(\sigma+1)p^2+(\sigma-1)q^2-6pq=(6+2\sigma)(p-p_*)^2$ with
$m^2=(1-\theta)pq$ proves the binary term. Stronger signed/mean/outer-gap
refinements retain their full ordinary proofs in~\cite[lines2470--4039]{source};
$D\ge0$ alone would not establish them.

If $n\ge N>0$, all costs are nonnegative and
\[
 (n/N)D_{\rm actual}\le d,\qquad
 D_{\rm actual}\ge b+gd_s+(N/n)c_\lambda\lambda,
\]
then $d\ge(n/N)(b+gd_s)+c_\lambda\lambda\ge b+gd_s+c_\lambda\lambda$.
This proves the basic weighted transfer. A later smaller marginal noise
ceiling does not replace the original joint intercept. Negative residuals
exclude; a zero residual forces nonnegative coordinates to zero.

\section{Normalization across one half}\label{app:norm}
The original sufficient binomial recipe required a Taylor segment below
$1/2$. The active recipe permits crossing under the following exact
extension, independently checked in the final audit.

\begin{lemma}[Symmetric derivative envelopes]
If $0<a<1/2$, $a<b<1$, $a+b<1$, then for $r\in[a,b]$,
\[
 r(1-r)\ge a(1-a),\qquad
 \frac{|1-2r|}{[r(1-r)]^j}\le\frac{1-2a}{[a(1-a)]^j}\quad(j=1,2).
\]
\end{lemma}
\begin{proof}
$[a,b]\subset[a,1-a]$. Set $s=\min(r,1-r)\in[a,1/2]$. The expressions
are symmetric; $s(1-s)$ increases and $1-2s$ decreases nonnegatively.
Their positive quotients therefore have their maxima at $a$.
\end{proof}
Let $n\ge N\ge1$, $p\in[p_-,p_+]\subset(0,1/2)$, $q=1-p$,
$1\le k\le n-1$, $\delta=k-np$, $r=k/n$, $v=pq$. Suppose the whole
segment between $p,r$ lies in $[a,b]$ above. If
\[
 E\ge\frac{1-2a}{a^2(1-a)^2},\quad L\ge\frac{1-2a}{a(1-a)},
 \quad R\ge\frac1{12a(1-a)},
\]
then the binomial mass $\pi_k$ satisfies
\[
 \sqrt n\,\pi_k=\frac{\phi(0)}{\sqrt v}
 e^{-\delta^2/(2nv)+\zeta},\qquad
 n|\zeta|<L|\delta|/2+R+E|\delta|^3/(6N).
\]
For proof, relative entropy has $D(p\Vert p)=D'(p\Vert p)=0$,
$D''(p\Vert p)=1/v$ and
$D^{(3)}(r\Vert p)=(2r-1)/[r^2(1-r)^2]$.
Taylor bounds $|nD(r\Vert p)-\delta^2/(2nv)|$ by
$E|\delta|^3/(6n^2)$; the logarithmic variance prefactor changes by at
most $L|\delta|/n$. Finite Stirling, $0<s_j<1/(12j)$, gives exactly
\[
 \sqrt n\,\pi_k=\frac{\phi(0)}{\sqrt{r(1-r)}}
 e^{-nD(r\Vert p)+s_n-s_k-s_{n-k}}.
\]
Writing $e=s_n-s_k-s_{n-k}$, use the compatible one-sided estimates
$e<1/(12n)$ and
$e>-1/(12k)-1/(12(n-k))=-1/[12nr(1-r)]$.
Thus $|e|<R/n$, and collecting all three errors proves the claim.
One must not add three incompatible absolute factorial estimates.

The literal code gates take $a=\texttt{left}$, $b=\texttt{right}$,
$a<p_- -\Delta/N$, $p_++\Delta/N<b$, $a+b<1$.
For $|\delta|\le\Delta$ they contain the whole segment. Gates
$Np_- -\Delta-2>0$ and $N(1-p_+)-\Delta-2>0$ ensure interior/adjacent
factorial domains. Positive $v_f<a(1-a)$ and upward-rounded
$(1-2a)/v_f^2$, $(1-2a)/v_f$, $1/(12v_f)$ supply $E,L,R$.
Mass bounds use $e^{-x}\ge1-x$, $e^x\le1/(1-x)$ only on their stated
positive domains. The original $p$ band remains below $1/2$ when bounding
$pq$; monotonicity is not assumed on the enlarged Taylor segment.
This lemma grants no physical acquisition or numerical exclusion.

\section{Exact source and audit map}
The ordinary source supplement~\cite{source} has SHA-256
\begin{center}\footnotesize\ttfamily
2b5c9419a5fc2363de21e02e200529b735\\c52e492f0d398aca6f42650c65eea1.
\end{center}
Exact applicable sections are attainment/contact368--463;
Fourier/resource1031--4039; physical COUNT4043--4765; small-index4822--6623;
and smoothing7458--7506 plus confinement/local7647--8887. These are ordinary theorem/proof
references, not numerical acceptance authorities. The longer calculations are provided in \path{analytic-supplement.txt},
whose opening index maps their lemmas to the claims used in this paper.

The completed review is
\path{work/final-audit/INDEPENDENT-MATHEMATICAL-AUDIT-v1.md}, incorporating
arithmetic continuations v1--v5, the resource-refinement review and the native
HIGH/support/barrier/confinement review. The actual global upper/lower/equality
objects have SHA prefixes \texttt{5acc019642ae}, \texttt{018bdf1b6918},
\texttt{8be2f7c6083d}; the full source/witness manifest supplies their identities.
Hashes locate objects; the arguments and complete predicates give meaning.
The companion supplement gives the ordinary proofs and their exact source scopes.

\begin{thebibliography}{9}
\bibitem{source}
K.~Rodziewicz and B.~Ko{\l}odziejek,
\emph{The iid Berry--Esseen constant: Analytic reductions and finite
verification}, working analytical source manuscript, October2026.
Selected ordinary proof sections are reproduced in the accompanying
\path{analytic-supplement.txt}, with their original labels and a dependency index.
\bibitem{schulz}
J.~Schulz, \emph{The optimal Berry--Esseen constant in the binomial case},
Dissertation, Universit\"at Trier,2016. Theorem1 and Remark4.1,
printed p.~15 (PDF p.~25).
\bibitem{shevtsova2013}
I.~Shevtsova, \emph{On the absolute constants in the Berry--Esseen inequality
and its structural and nonuniform improvements},
Informatics and Applications7(1) (2013),124--125 (in Russian).
\bibitem{shevtsova2012}
I.~Shevtsova, \emph{Moment-type estimates with asymptotically optimal
structure for the accuracy of the normal approximation},
Annales Mathematicae et Informaticae39 (2012),241--307;
Lemma4.1, printed p.~263.
\end{thebibliography}
\end{document}
